\begin{beginnerstatement}[Kurz erklärt: Einstieg für Entwickler]

\emph{Onboarding} ist der Weg, auf dem ein Entwickler das Projekt und seine
Arbeitsabläufe erstmals einrichtet. Ein Repository zu klonen bedeutet, mit Git
eine lokale Arbeitskopie davon anzulegen. Die Entwicklerumgebung umfasst die
benötigten Programme, Abhängigkeiten und Einstellungen auf dem Arbeitsrechner.
Das Master-Repository ist die gemeinsam gepflegte Vorlage; der eigentliche
Produktcode gehört danach in das getrennte Produktrepository. Die Reihenfolge
lautet: Master klonen \(\rightarrow\) \texttt{doctor} \(\rightarrow\)
\texttt{init} \(\rightarrow\) Zielprojekt \(\rightarrow\) erneut
\texttt{doctor} \(\rightarrow\) \texttt{install} \(\rightarrow\)
\texttt{run}. Damit wird erst die Vorlage geprüft und anschließend das
erzeugte Zielprojekt eingerichtet und gestartet.

\end{beginnerstatement}
